% s3_calibration.tex — Calibration section
% Paper: "Tokenized Housing and Lifecycle Portfolio Choice:
%         A Decoupling of Location from Housing Exposure"
%
% Source: docs/calibration_v3.md (cloud agent fire 5)
%         docs/welfare_decomp_v4.md (channel decomposition plan)
%         paper/outline_v4.md (section structure)
%
% Status: complete draft; one parameter (rho_AB sensitivity) awaits
%         server1 baseline results before final anchor confirmation (H2' gate).
%
% Compilation: included from main.tex via % s3_calibration.tex — Calibration section
% Paper: "Tokenized Housing and Lifecycle Portfolio Choice:
%         A Decoupling of Location from Housing Exposure"
%
% Source: docs/calibration_v3.md (cloud agent fire 5)
%         docs/welfare_decomp_v4.md (channel decomposition plan)
%         paper/outline_v4.md (section structure)
%
% Status: complete draft; one parameter (rho_AB sensitivity) awaits
%         server1 baseline results before final anchor confirmation (H2' gate).
%
% Compilation: included from main.tex via % s3_calibration.tex — Calibration section
% Paper: "Tokenized Housing and Lifecycle Portfolio Choice:
%         A Decoupling of Location from Housing Exposure"
%
% Source: docs/calibration_v3.md (cloud agent fire 5)
%         docs/welfare_decomp_v4.md (channel decomposition plan)
%         paper/outline_v4.md (section structure)
%
% Status: complete draft; one parameter (rho_AB sensitivity) awaits
%         server1 baseline results before final anchor confirmation (H2' gate).
%
% Compilation: included from main.tex via % s3_calibration.tex — Calibration section
% Paper: "Tokenized Housing and Lifecycle Portfolio Choice:
%         A Decoupling of Location from Housing Exposure"
%
% Source: docs/calibration_v3.md (cloud agent fire 5)
%         docs/welfare_decomp_v4.md (channel decomposition plan)
%         paper/outline_v4.md (section structure)
%
% Status: complete draft; one parameter (rho_AB sensitivity) awaits
%         server1 baseline results before final anchor confirmation (H2' gate).
%
% Compilation: included from main.tex via \input{sections/s3_calibration}
%              Requires amsmath, booktabs, multirow, siunitx (for tabular).

% ─────────────────────────────────────────────────────────────────────────────
\section{Calibration}
\label{sec:calibration}
% ─────────────────────────────────────────────────────────────────────────────

Table~\ref{tab:parameters} summarises the baseline parameter values.
The calibration follows the household-finance lifecycle tradition
(Cocco, Gomes, and Maenhout, 2005; Yao and Zhang, 2005), extended
with three v4-specific objects: the relocation probability $p_{\text{work}}$,
the round-trip transaction-cost pair $(\tau_{\text{sell}}, \tau_{\text{buy}})$,
and the cross-location idiosyncratic correlation $\rho_{AB}$.
We discuss each group in turn.

% ── Table 1: Parameter summary ───────────────────────────────────────────────
\begin{table}[t]
\centering
\caption{Baseline Parameters}
\label{tab:parameters}
\small
\begin{tabular}{llcl}
\toprule
Parameter & Symbol & Value & Source \\
\midrule
\multicolumn{4}{l}{\textit{Preferences and lifecycle}} \\
Risk aversion         & $\gamma$         & 5.0  & CGM (2005) \\
Time discount         & $\beta$          & 0.96 & CGM (2005) \\
Entry age / retire / terminal & --- & 25 / 65 / 80 & Standard \\
\addlinespace
\multicolumn{4}{l}{\textit{Financial assets}} \\
Risk-free gross return     & $R_f$           & 1.02  & CGM (2005) \\
Equity risk premium        & $\mu_S - r_f$   & 4\%   & CGM (2005) \\
Stock return volatility    & $\sigma_S$      & 15.7\% & CGM (2005) \\
\addlinespace
\multicolumn{4}{l}{\textit{Income process (CGM 2005)}} \\
Permanent shock variance   & $\sigma_u^2$    & 0.0106 & CGM Table~I \\
Transitory shock variance  & $\sigma_\varepsilon^2$ & 0.0738 & CGM Table~I \\
Retirement replacement     & $\lambda$       & 0.65  & CGM (2005) \\
\addlinespace
\multicolumn{4}{l}{\textit{Housing return (Cocco 2005; Yao-Zhang 2005)}} \\
Expected log price growth  & $g_h$           & 1.6\%  & Cocco (2005) \\
Total housing return std   & $\sigma_h$      & 11.5\% & Cocco (2005) \\
Aggregate factor std       & $\sigma_{\text{div}}$ & 10.0\% & Case-Shiller / NAREIT \\
Idiosyncratic std (derived) & $\sigma_\iota$ & 5.7\%  & $\sqrt{\sigma_h^2 - \sigma_{\text{div}}^2}$ \\
Rent-to-price ratio        & $\varrho$       & 5.0\%  & Yao and Zhang (2005) \\
Maintenance-to-price ratio & $m$             & 1.0\%  & Cocco (2005) \\
Rent-saving wedge (derived) & $\delta_{\text{own}} = \varrho - m$ & 4.0\% & --- \\
\addlinespace
\multicolumn{4}{l}{\textit{Mobility (PSID-anchored)}} \\
Relocation prob, working age & $p_{\text{work}}$ & 6.0\% & PSID inter-MSA \\
Relocation prob, retired     & $p_{\text{ret}}$  & 2.0\% & PSID retired \\
\addlinespace
\multicolumn{4}{l}{\textit{Transaction costs}} \\
Traditional sell cost  & $\tau_{\text{sell}}$ & 6.0\% & NAR (2023) \\
Traditional buy cost   & $\tau_{\text{buy}}$  & 2.5\% & CFPB / ALTA \\
Token transfer cost    & $\tau_{\text{token}}$& 1.0\% & RealT / Lofty surveys \\
\addlinespace
\multicolumn{4}{l}{\textit{Cross-location correlation}} \\
Idiosyncratic corr., MSA pair & $\rho_{AB}$ & 0.50 & Case-Shiller (2000--2023) \\
\addlinespace
\multicolumn{4}{l}{\textit{Mortgage (baseline off)}} \\
LTV constraint & $\overline{\text{ltv}}$ & 0.0 & Baseline no-mortgage \\
\bottomrule
\end{tabular}
\end{table}

\subsection{Preferences and Income Process}
\label{sec:calibration:prefs}

Preferences and the income process are taken directly from
Cocco, Gomes, and Maenhout (2005, CGM hereafter).
Risk aversion $\gamma = 5$ and the discount factor $\beta = 0.96$ are
their baseline values.
The labor income age profile follows the CGM cubic polynomial fit
to PSID data:
\begin{equation}
    f(\text{age}) \;=\;
    -2.170 + 0.168\,(\text{age}/10)
    - 0.032\,(\text{age}/10)^2
    + 0.002\,(\text{age}/10)^3.
    \label{eq:income_profile}
\end{equation}
Permanent and transitory income shocks have variances $\sigma_u^2 = 0.0106$
and $\sigma_\varepsilon^2 = 0.0738$ (CGM Table~I).
The retirement replacement rate $\lambda = 0.65$ matches the median
Social Security replacement rate at median lifetime earnings.

Financial asset parameters follow the same source:
the equity premium is $4\%$, stock return volatility $15.7\%$,
and the risk-free rate $2\%$ real.

\subsection{Housing Return Decomposition}
\label{sec:calibration:housing}

Single-location housing return volatility $\sigma_h = 11.5\%$ is from
Cocco (2005).
We decompose the single-location log return into a shared aggregate factor
$\eta_{\text{div}}$ and an idiosyncratic component $\iota_\ell$:
\begin{equation}
    \log R_\ell \;=\; \mu_h + \eta_{\text{div}} + \iota_\ell,
    \qquad
    \eta_{\text{div}} \sim \mathcal{N}(0,\sigma_{\text{div}}^2),
    \quad
    \iota_\ell \sim \mathcal{N}(0,\sigma_\iota^2),
    \label{eq:return_decomp}
\end{equation}
with $\sigma_h^2 = \sigma_{\text{div}}^2 + \sigma_\iota^2$.
We set $\sigma_{\text{div}} = 10\%$ so that the aggregate factor accounts
for approximately $75\%$ of total single-location housing variance
($\sigma_{\text{div}}^2 / \sigma_h^2 \approx 0.756$), consistent with
the common factor dominating MSA-level variance in Case-Shiller data
and with Piazzesi and Schneider~(2016).
The idiosyncratic component $\sigma_\iota = \sqrt{0.115^2 - 0.10^2} \approx 5.7\%$
is calibrated residually.

The idiosyncratic components at the two locations follow a bivariate normal
with correlation $\rho_{AB}$ (Section~\ref{sec:calibration:rhoAB}).
The Cholesky factorisation is:
\begin{align}
    \iota_A &= \sigma_\iota\,\xi_A, \label{eq:chol_A} \\
    \iota_B &= \rho_{AB}\,\iota_A + \sqrt{1 - \rho_{AB}^2}\,\sigma_\iota\,\xi_B,
    \label{eq:chol_B}
\end{align}
with $\xi_A, \xi_B \sim \mathcal{N}(0,1)$ independent.

The rent-to-price ratio $\varrho = 5\%$ follows Yao and Zhang~(2005).
Maintenance costs $m = 1\%$ follow Cocco~(2005).
The rent-saving wedge $\delta_{\text{own}} = \varrho - m = 4\%$ governs
the per-period benefit of occupying a fractional token position
($\kappa$ rule, equation~(\ref{eq:kappa}) in Section~\ref{sec:model:kappa}).

\subsection{Mobility Rate}
\label{sec:calibration:mobility}

We anchor the relocation probability to PSID inter-MSA mobility data.
Total annual relocation rates in the PSID average $7$--$12\%$ for
working-age households but include short-distance, within-MSA moves that
do not change the relevant housing market.
Long-distance, inter-MSA moves account for approximately $40$--$50\%$
of total PSID mobility, yielding an inter-MSA rate of $5$--$7\%$ annually
for prime-age workers.
We set $p_{\text{work}} = 6\%$ (PSID midpoint) for working ages
$25$--$64$ and $p_{\text{ret}} = 2\%$ for ages $65$--$80$.

This choice is consistent with Bagliano, Fugazza, and Nicodano~(2014, RFS),
who calibrate career-driven mobility shocks at $5$--$8\%$ annually, and
broadly in the range of Yao and Zhang~(2005) ($4\%$) and Saks and
Wozniak~(2011, JLE) ($2$--$4\%$ inter-state, implying $4$--$8\%$ inter-MSA).

The sensitivity grid (Table~\ref{tab:sensitivity_grid} below)
spans $p_{\text{work}} \in \{0, 0.02, 0.06, 0.12\}$.
At $p_{\text{work}} = 0$, the cross-location hedge channel must vanish by
construction (no relocation, no cost-saving motive); this is a pre-registered
falsification test (Referee Round~4, test~r).

\subsection{Transaction Costs}
\label{sec:calibration:txcost}

\paragraph{Traditional housing ($E1\_2L$).}
We set the seller commission and closing costs at $\tau_{\text{sell}} = 6\%$,
consistent with the NAR~(2023) aggregate for a typical single-family
home sale ($5$--$6\%$ agent commission plus $0.5$--$1\%$ closing costs).
The buyer cost is $\tau_{\text{buy}} = 2.5\%$,
covering origination fees ($0.75\%$), title insurance and escrow
($0.75\%$), appraisal and inspection ($0.25\%$), and prepaid items
($0.25\%$), consistent with CFPB and ALTA surveys.
The implied round-trip cost for an $E1\_2L$ relocation is
$\tau_{\text{sell}} + \tau_{\text{buy}} = 8.5\%$,
within the $5$--$10\%$ range used in the structural housing literature
(Diaz and Luengo-Prado, 2008; Head, Lloyd-Ellis, and Sun, 2014).

\paragraph{Tokenized housing ($E2\_2L$).}
Platform surveys of RealT, Lofty, and DigiShares~(2024) report trading
fees in the range $0.1$--$0.5\%$ plus blockchain gas and settlement
of $0.05$--$0.3\%$ on Layer-2 networks.
We set $\tau_{\text{token}} = 1\%$ to include a regulatory-compliance
margin of $0.5\%$ (reflecting SEC Regulation~D and AML overhead).
The structural distinction from $E1\_2L$ is that $\tau_{\text{token}} \ll \tau_{\text{sell}} + \tau_{\text{buy}}$:
tokens are portable across relocations, so the $8.5\%$ round-trip is avoided
for every unit of cross-location position retained.

\paragraph{Expected hedge premium.}
The per-period value of pre-positioning one unit of $x_B$ while at
location $A$ is
\begin{equation}
    \text{hedge premium per unit}
    \;=\;
    p_{\text{work}} \times \tau_{\text{buy}}
    \;\approx\;
    0.06 \times 0.025
    \;=\; 0.15\%\ \text{per period}.
    \label{eq:hedge_premium}
\end{equation}
This is the incremental saving on buying costs that accrues when the
household has pre-positioned $x_B > 0$ prior to relocation.
Accumulated over a working-life horizon, equation~(\ref{eq:hedge_premium})
implies a lifetime CEV contribution of $1$--$2\%$ from the hedge
channel alone, with the continuous-$x$ rent-saving channel adding a further
$3$--$4\%$ (see Section~\ref{sec:results:decomp}).

\subsection{Cross-Location House-Price Correlation}
\label{sec:calibration:rhoAB}

The parameter $\rho_{AB}$ governs the idiosyncratic correlation between
location $A$ and location $B$ returns.
Using the decomposition~(\ref{eq:return_decomp}),
the observed return correlation between two MSAs is:
\begin{equation}
    \operatorname{Corr}(R_A, R_B)
    \;\approx\;
    \frac{\sigma_{\text{div}}^2}{\sigma_h^2}
    \;+\;
    \frac{\sigma_\iota^2}{\sigma_h^2}\,\rho_{AB}
    \;\approx\; 0.756 + 0.244\,\rho_{AB}.
    \label{eq:rhoAB_mapping}
\end{equation}
We set $\rho_{AB} = 0.50$, which maps to an observed return correlation of
$\approx 0.88$, consistent with pairs of geographically proximate
US metros in the Case-Shiller 20-city panel (2000--2023).
This is the midpoint of the $0.30$--$0.70$ range that covers inter-regional
moves (different metro areas): at one extreme, Chicago--Phoenix
($\hat{\rho}_{AB} \approx 0.35$) and at the other,
NY--Boston ($\hat{\rho}_{AB} \approx 0.60$).

Equation~(\ref{eq:rhoAB_mapping}) implies that the hedge value of
$x_B$ at location $A$ derives from bearing the idiosyncratic component
$\iota_B$ independently of $\iota_A$.
As $\rho_{AB} \to 1$, locations $A$ and $B$ become nearly identical
markets and $x_B$ is a near-substitute for $x_A$; the hedge channel
collapses.
The Round-4 referee pre-registered test~(m) requires $\operatorname{CEV}(E2\_2L \text{ vs } E1\_2L)$
to decrease significantly at $\rho_{AB} = 0.95$.

\subsection{Sensitivity Grid}
\label{sec:calibration:sensitivity}

Table~\ref{tab:sensitivity_grid} lists the pre-registered sensitivity
sweeps queued for execution after the v4 baseline results are available.
The three P1-priority sweeps cover the three parameters most directly
linked to the welfare decomposition channels.

\begin{table}[t]
\centering
\caption{Pre-Registered Sensitivity Sweeps}
\label{tab:sensitivity_grid}
\small
\begin{tabular}{llll}
\toprule
Parameter & Sweep values & Channel & Script \\
\midrule
$\rho_{AB}$          & $\{0,\, 0.25,\, 0.50,\, 0.75,\, 0.95\}$  & Hedge (maintained) & \texttt{sweep\_rhoAB.sh} \\
$p_{\text{work}}$    & $\{0,\, 0.02,\, 0.06,\, 0.12\}$           & Hedge (frequency)  & \texttt{sweep\_prelocate.sh} \\
$\tau_{\text{buy}}$  & $\{0,\, 0.025,\, 0.04,\, 0.06\}$          & Avoided tx-cost    & \texttt{sweep\_txcost.sh} \\
\midrule
\multicolumn{4}{l}{\textit{Phase 2 sweeps (after H2' approval)}} \\
$\gamma$             & $\{3,\, 5,\, 8\}$                         & Risk channel       & --- \\
$\delta_{\text{own}}$& $\{0.02,\, 0.04,\, 0.06\}$               & Rent-saving        & --- \\
$\sigma_{\text{div}}$& $\{0.08,\, 0.10,\, 0.12\}$               & Idio.\ risk share  & --- \\
\bottomrule
\end{tabular}
\end{table}

At $p_{\text{work}} = 0$ and at $\tau_{\text{buy}} = 0$,
the hedge premium~(\ref{eq:hedge_premium}) is exactly zero; any
remaining welfare gap between $E2\_2L$ and $E1\_2L$ in these cells
reflects the continuous-$x$ rent-saving channel only and serves
as a falsification bound.
At $\rho_{AB} = 0.95$, the maintained-hedge channel must also collapse
because $x_B$ and $x_A$ are near-substitutes.
These three boundary conditions define sharp predictions that distinguish
the two decomposed channels (Section~\ref{sec:results:decomp}).

%              Requires amsmath, booktabs, multirow, siunitx (for tabular).

% ─────────────────────────────────────────────────────────────────────────────
\section{Calibration}
\label{sec:calibration}
% ─────────────────────────────────────────────────────────────────────────────

Table~\ref{tab:parameters} summarises the baseline parameter values.
The calibration follows the household-finance lifecycle tradition
(Cocco, Gomes, and Maenhout, 2005; Yao and Zhang, 2005), extended
with three v4-specific objects: the relocation probability $p_{\text{work}}$,
the round-trip transaction-cost pair $(\tau_{\text{sell}}, \tau_{\text{buy}})$,
and the cross-location idiosyncratic correlation $\rho_{AB}$.
We discuss each group in turn.

% ── Table 1: Parameter summary ───────────────────────────────────────────────
\begin{table}[t]
\centering
\caption{Baseline Parameters}
\label{tab:parameters}
\small
\begin{tabular}{llcl}
\toprule
Parameter & Symbol & Value & Source \\
\midrule
\multicolumn{4}{l}{\textit{Preferences and lifecycle}} \\
Risk aversion         & $\gamma$         & 5.0  & CGM (2005) \\
Time discount         & $\beta$          & 0.96 & CGM (2005) \\
Entry age / retire / terminal & --- & 25 / 65 / 80 & Standard \\
\addlinespace
\multicolumn{4}{l}{\textit{Financial assets}} \\
Risk-free gross return     & $R_f$           & 1.02  & CGM (2005) \\
Equity risk premium        & $\mu_S - r_f$   & 4\%   & CGM (2005) \\
Stock return volatility    & $\sigma_S$      & 15.7\% & CGM (2005) \\
\addlinespace
\multicolumn{4}{l}{\textit{Income process (CGM 2005)}} \\
Permanent shock variance   & $\sigma_u^2$    & 0.0106 & CGM Table~I \\
Transitory shock variance  & $\sigma_\varepsilon^2$ & 0.0738 & CGM Table~I \\
Retirement replacement     & $\lambda$       & 0.65  & CGM (2005) \\
\addlinespace
\multicolumn{4}{l}{\textit{Housing return (Cocco 2005; Yao-Zhang 2005)}} \\
Expected log price growth  & $g_h$           & 1.6\%  & Cocco (2005) \\
Total housing return std   & $\sigma_h$      & 11.5\% & Cocco (2005) \\
Aggregate factor std       & $\sigma_{\text{div}}$ & 10.0\% & Case-Shiller / NAREIT \\
Idiosyncratic std (derived) & $\sigma_\iota$ & 5.7\%  & $\sqrt{\sigma_h^2 - \sigma_{\text{div}}^2}$ \\
Rent-to-price ratio        & $\varrho$       & 5.0\%  & Yao and Zhang (2005) \\
Maintenance-to-price ratio & $m$             & 1.0\%  & Cocco (2005) \\
Rent-saving wedge (derived) & $\delta_{\text{own}} = \varrho - m$ & 4.0\% & --- \\
\addlinespace
\multicolumn{4}{l}{\textit{Mobility (PSID-anchored)}} \\
Relocation prob, working age & $p_{\text{work}}$ & 6.0\% & PSID inter-MSA \\
Relocation prob, retired     & $p_{\text{ret}}$  & 2.0\% & PSID retired \\
\addlinespace
\multicolumn{4}{l}{\textit{Transaction costs}} \\
Traditional sell cost  & $\tau_{\text{sell}}$ & 6.0\% & NAR (2023) \\
Traditional buy cost   & $\tau_{\text{buy}}$  & 2.5\% & CFPB / ALTA \\
Token transfer cost    & $\tau_{\text{token}}$& 1.0\% & RealT / Lofty surveys \\
\addlinespace
\multicolumn{4}{l}{\textit{Cross-location correlation}} \\
Idiosyncratic corr., MSA pair & $\rho_{AB}$ & 0.50 & Case-Shiller (2000--2023) \\
\addlinespace
\multicolumn{4}{l}{\textit{Mortgage (baseline off)}} \\
LTV constraint & $\overline{\text{ltv}}$ & 0.0 & Baseline no-mortgage \\
\bottomrule
\end{tabular}
\end{table}

\subsection{Preferences and Income Process}
\label{sec:calibration:prefs}

Preferences and the income process are taken directly from
Cocco, Gomes, and Maenhout (2005, CGM hereafter).
Risk aversion $\gamma = 5$ and the discount factor $\beta = 0.96$ are
their baseline values.
The labor income age profile follows the CGM cubic polynomial fit
to PSID data:
\begin{equation}
    f(\text{age}) \;=\;
    -2.170 + 0.168\,(\text{age}/10)
    - 0.032\,(\text{age}/10)^2
    + 0.002\,(\text{age}/10)^3.
    \label{eq:income_profile}
\end{equation}
Permanent and transitory income shocks have variances $\sigma_u^2 = 0.0106$
and $\sigma_\varepsilon^2 = 0.0738$ (CGM Table~I).
The retirement replacement rate $\lambda = 0.65$ matches the median
Social Security replacement rate at median lifetime earnings.

Financial asset parameters follow the same source:
the equity premium is $4\%$, stock return volatility $15.7\%$,
and the risk-free rate $2\%$ real.

\subsection{Housing Return Decomposition}
\label{sec:calibration:housing}

Single-location housing return volatility $\sigma_h = 11.5\%$ is from
Cocco (2005).
We decompose the single-location log return into a shared aggregate factor
$\eta_{\text{div}}$ and an idiosyncratic component $\iota_\ell$:
\begin{equation}
    \log R_\ell \;=\; \mu_h + \eta_{\text{div}} + \iota_\ell,
    \qquad
    \eta_{\text{div}} \sim \mathcal{N}(0,\sigma_{\text{div}}^2),
    \quad
    \iota_\ell \sim \mathcal{N}(0,\sigma_\iota^2),
    \label{eq:return_decomp}
\end{equation}
with $\sigma_h^2 = \sigma_{\text{div}}^2 + \sigma_\iota^2$.
We set $\sigma_{\text{div}} = 10\%$ so that the aggregate factor accounts
for approximately $75\%$ of total single-location housing variance
($\sigma_{\text{div}}^2 / \sigma_h^2 \approx 0.756$), consistent with
the common factor dominating MSA-level variance in Case-Shiller data
and with Piazzesi and Schneider~(2016).
The idiosyncratic component $\sigma_\iota = \sqrt{0.115^2 - 0.10^2} \approx 5.7\%$
is calibrated residually.

The idiosyncratic components at the two locations follow a bivariate normal
with correlation $\rho_{AB}$ (Section~\ref{sec:calibration:rhoAB}).
The Cholesky factorisation is:
\begin{align}
    \iota_A &= \sigma_\iota\,\xi_A, \label{eq:chol_A} \\
    \iota_B &= \rho_{AB}\,\iota_A + \sqrt{1 - \rho_{AB}^2}\,\sigma_\iota\,\xi_B,
    \label{eq:chol_B}
\end{align}
with $\xi_A, \xi_B \sim \mathcal{N}(0,1)$ independent.

The rent-to-price ratio $\varrho = 5\%$ follows Yao and Zhang~(2005).
Maintenance costs $m = 1\%$ follow Cocco~(2005).
The rent-saving wedge $\delta_{\text{own}} = \varrho - m = 4\%$ governs
the per-period benefit of occupying a fractional token position
($\kappa$ rule, equation~(\ref{eq:kappa}) in Section~\ref{sec:model:kappa}).

\subsection{Mobility Rate}
\label{sec:calibration:mobility}

We anchor the relocation probability to PSID inter-MSA mobility data.
Total annual relocation rates in the PSID average $7$--$12\%$ for
working-age households but include short-distance, within-MSA moves that
do not change the relevant housing market.
Long-distance, inter-MSA moves account for approximately $40$--$50\%$
of total PSID mobility, yielding an inter-MSA rate of $5$--$7\%$ annually
for prime-age workers.
We set $p_{\text{work}} = 6\%$ (PSID midpoint) for working ages
$25$--$64$ and $p_{\text{ret}} = 2\%$ for ages $65$--$80$.

This choice is consistent with Bagliano, Fugazza, and Nicodano~(2014, RFS),
who calibrate career-driven mobility shocks at $5$--$8\%$ annually, and
broadly in the range of Yao and Zhang~(2005) ($4\%$) and Saks and
Wozniak~(2011, JLE) ($2$--$4\%$ inter-state, implying $4$--$8\%$ inter-MSA).

The sensitivity grid (Table~\ref{tab:sensitivity_grid} below)
spans $p_{\text{work}} \in \{0, 0.02, 0.06, 0.12\}$.
At $p_{\text{work}} = 0$, the cross-location hedge channel must vanish by
construction (no relocation, no cost-saving motive); this is a pre-registered
falsification test (Referee Round~4, test~r).

\subsection{Transaction Costs}
\label{sec:calibration:txcost}

\paragraph{Traditional housing ($E1\_2L$).}
We set the seller commission and closing costs at $\tau_{\text{sell}} = 6\%$,
consistent with the NAR~(2023) aggregate for a typical single-family
home sale ($5$--$6\%$ agent commission plus $0.5$--$1\%$ closing costs).
The buyer cost is $\tau_{\text{buy}} = 2.5\%$,
covering origination fees ($0.75\%$), title insurance and escrow
($0.75\%$), appraisal and inspection ($0.25\%$), and prepaid items
($0.25\%$), consistent with CFPB and ALTA surveys.
The implied round-trip cost for an $E1\_2L$ relocation is
$\tau_{\text{sell}} + \tau_{\text{buy}} = 8.5\%$,
within the $5$--$10\%$ range used in the structural housing literature
(Diaz and Luengo-Prado, 2008; Head, Lloyd-Ellis, and Sun, 2014).

\paragraph{Tokenized housing ($E2\_2L$).}
Platform surveys of RealT, Lofty, and DigiShares~(2024) report trading
fees in the range $0.1$--$0.5\%$ plus blockchain gas and settlement
of $0.05$--$0.3\%$ on Layer-2 networks.
We set $\tau_{\text{token}} = 1\%$ to include a regulatory-compliance
margin of $0.5\%$ (reflecting SEC Regulation~D and AML overhead).
The structural distinction from $E1\_2L$ is that $\tau_{\text{token}} \ll \tau_{\text{sell}} + \tau_{\text{buy}}$:
tokens are portable across relocations, so the $8.5\%$ round-trip is avoided
for every unit of cross-location position retained.

\paragraph{Expected hedge premium.}
The per-period value of pre-positioning one unit of $x_B$ while at
location $A$ is
\begin{equation}
    \text{hedge premium per unit}
    \;=\;
    p_{\text{work}} \times \tau_{\text{buy}}
    \;\approx\;
    0.06 \times 0.025
    \;=\; 0.15\%\ \text{per period}.
    \label{eq:hedge_premium}
\end{equation}
This is the incremental saving on buying costs that accrues when the
household has pre-positioned $x_B > 0$ prior to relocation.
Accumulated over a working-life horizon, equation~(\ref{eq:hedge_premium})
implies a lifetime CEV contribution of $1$--$2\%$ from the hedge
channel alone, with the continuous-$x$ rent-saving channel adding a further
$3$--$4\%$ (see Section~\ref{sec:results:decomp}).

\subsection{Cross-Location House-Price Correlation}
\label{sec:calibration:rhoAB}

The parameter $\rho_{AB}$ governs the idiosyncratic correlation between
location $A$ and location $B$ returns.
Using the decomposition~(\ref{eq:return_decomp}),
the observed return correlation between two MSAs is:
\begin{equation}
    \operatorname{Corr}(R_A, R_B)
    \;\approx\;
    \frac{\sigma_{\text{div}}^2}{\sigma_h^2}
    \;+\;
    \frac{\sigma_\iota^2}{\sigma_h^2}\,\rho_{AB}
    \;\approx\; 0.756 + 0.244\,\rho_{AB}.
    \label{eq:rhoAB_mapping}
\end{equation}
We set $\rho_{AB} = 0.50$, which maps to an observed return correlation of
$\approx 0.88$, consistent with pairs of geographically proximate
US metros in the Case-Shiller 20-city panel (2000--2023).
This is the midpoint of the $0.30$--$0.70$ range that covers inter-regional
moves (different metro areas): at one extreme, Chicago--Phoenix
($\hat{\rho}_{AB} \approx 0.35$) and at the other,
NY--Boston ($\hat{\rho}_{AB} \approx 0.60$).

Equation~(\ref{eq:rhoAB_mapping}) implies that the hedge value of
$x_B$ at location $A$ derives from bearing the idiosyncratic component
$\iota_B$ independently of $\iota_A$.
As $\rho_{AB} \to 1$, locations $A$ and $B$ become nearly identical
markets and $x_B$ is a near-substitute for $x_A$; the hedge channel
collapses.
The Round-4 referee pre-registered test~(m) requires $\operatorname{CEV}(E2\_2L \text{ vs } E1\_2L)$
to decrease significantly at $\rho_{AB} = 0.95$.

\subsection{Sensitivity Grid}
\label{sec:calibration:sensitivity}

Table~\ref{tab:sensitivity_grid} lists the pre-registered sensitivity
sweeps queued for execution after the v4 baseline results are available.
The three P1-priority sweeps cover the three parameters most directly
linked to the welfare decomposition channels.

\begin{table}[t]
\centering
\caption{Pre-Registered Sensitivity Sweeps}
\label{tab:sensitivity_grid}
\small
\begin{tabular}{llll}
\toprule
Parameter & Sweep values & Channel & Script \\
\midrule
$\rho_{AB}$          & $\{0,\, 0.25,\, 0.50,\, 0.75,\, 0.95\}$  & Hedge (maintained) & \texttt{sweep\_rhoAB.sh} \\
$p_{\text{work}}$    & $\{0,\, 0.02,\, 0.06,\, 0.12\}$           & Hedge (frequency)  & \texttt{sweep\_prelocate.sh} \\
$\tau_{\text{buy}}$  & $\{0,\, 0.025,\, 0.04,\, 0.06\}$          & Avoided tx-cost    & \texttt{sweep\_txcost.sh} \\
\midrule
\multicolumn{4}{l}{\textit{Phase 2 sweeps (after H2' approval)}} \\
$\gamma$             & $\{3,\, 5,\, 8\}$                         & Risk channel       & --- \\
$\delta_{\text{own}}$& $\{0.02,\, 0.04,\, 0.06\}$               & Rent-saving        & --- \\
$\sigma_{\text{div}}$& $\{0.08,\, 0.10,\, 0.12\}$               & Idio.\ risk share  & --- \\
\bottomrule
\end{tabular}
\end{table}

At $p_{\text{work}} = 0$ and at $\tau_{\text{buy}} = 0$,
the hedge premium~(\ref{eq:hedge_premium}) is exactly zero; any
remaining welfare gap between $E2\_2L$ and $E1\_2L$ in these cells
reflects the continuous-$x$ rent-saving channel only and serves
as a falsification bound.
At $\rho_{AB} = 0.95$, the maintained-hedge channel must also collapse
because $x_B$ and $x_A$ are near-substitutes.
These three boundary conditions define sharp predictions that distinguish
the two decomposed channels (Section~\ref{sec:results:decomp}).

%              Requires amsmath, booktabs, multirow, siunitx (for tabular).

% ─────────────────────────────────────────────────────────────────────────────
\section{Calibration}
\label{sec:calibration}
% ─────────────────────────────────────────────────────────────────────────────

Table~\ref{tab:parameters} summarises the baseline parameter values.
The calibration follows the household-finance lifecycle tradition
(Cocco, Gomes, and Maenhout, 2005; Yao and Zhang, 2005), extended
with three v4-specific objects: the relocation probability $p_{\text{work}}$,
the round-trip transaction-cost pair $(\tau_{\text{sell}}, \tau_{\text{buy}})$,
and the cross-location idiosyncratic correlation $\rho_{AB}$.
We discuss each group in turn.

% ── Table 1: Parameter summary ───────────────────────────────────────────────
\begin{table}[t]
\centering
\caption{Baseline Parameters}
\label{tab:parameters}
\small
\begin{tabular}{llcl}
\toprule
Parameter & Symbol & Value & Source \\
\midrule
\multicolumn{4}{l}{\textit{Preferences and lifecycle}} \\
Risk aversion         & $\gamma$         & 5.0  & CGM (2005) \\
Time discount         & $\beta$          & 0.96 & CGM (2005) \\
Entry age / retire / terminal & --- & 25 / 65 / 80 & Standard \\
\addlinespace
\multicolumn{4}{l}{\textit{Financial assets}} \\
Risk-free gross return     & $R_f$           & 1.02  & CGM (2005) \\
Equity risk premium        & $\mu_S - r_f$   & 4\%   & CGM (2005) \\
Stock return volatility    & $\sigma_S$      & 15.7\% & CGM (2005) \\
\addlinespace
\multicolumn{4}{l}{\textit{Income process (CGM 2005)}} \\
Permanent shock variance   & $\sigma_u^2$    & 0.0106 & CGM Table~I \\
Transitory shock variance  & $\sigma_\varepsilon^2$ & 0.0738 & CGM Table~I \\
Retirement replacement     & $\lambda$       & 0.65  & CGM (2005) \\
\addlinespace
\multicolumn{4}{l}{\textit{Housing return (Cocco 2005; Yao-Zhang 2005)}} \\
Expected log price growth  & $g_h$           & 1.6\%  & Cocco (2005) \\
Total housing return std   & $\sigma_h$      & 11.5\% & Cocco (2005) \\
Aggregate factor std       & $\sigma_{\text{div}}$ & 10.0\% & Case-Shiller / NAREIT \\
Idiosyncratic std (derived) & $\sigma_\iota$ & 5.7\%  & $\sqrt{\sigma_h^2 - \sigma_{\text{div}}^2}$ \\
Rent-to-price ratio        & $\varrho$       & 5.0\%  & Yao and Zhang (2005) \\
Maintenance-to-price ratio & $m$             & 1.0\%  & Cocco (2005) \\
Rent-saving wedge (derived) & $\delta_{\text{own}} = \varrho - m$ & 4.0\% & --- \\
\addlinespace
\multicolumn{4}{l}{\textit{Mobility (PSID-anchored)}} \\
Relocation prob, working age & $p_{\text{work}}$ & 6.0\% & PSID inter-MSA \\
Relocation prob, retired     & $p_{\text{ret}}$  & 2.0\% & PSID retired \\
\addlinespace
\multicolumn{4}{l}{\textit{Transaction costs}} \\
Traditional sell cost  & $\tau_{\text{sell}}$ & 6.0\% & NAR (2023) \\
Traditional buy cost   & $\tau_{\text{buy}}$  & 2.5\% & CFPB / ALTA \\
Token transfer cost    & $\tau_{\text{token}}$& 1.0\% & RealT / Lofty surveys \\
\addlinespace
\multicolumn{4}{l}{\textit{Cross-location correlation}} \\
Idiosyncratic corr., MSA pair & $\rho_{AB}$ & 0.50 & Case-Shiller (2000--2023) \\
\addlinespace
\multicolumn{4}{l}{\textit{Mortgage (baseline off)}} \\
LTV constraint & $\overline{\text{ltv}}$ & 0.0 & Baseline no-mortgage \\
\bottomrule
\end{tabular}
\end{table}

\subsection{Preferences and Income Process}
\label{sec:calibration:prefs}

Preferences and the income process are taken directly from
Cocco, Gomes, and Maenhout (2005, CGM hereafter).
Risk aversion $\gamma = 5$ and the discount factor $\beta = 0.96$ are
their baseline values.
The labor income age profile follows the CGM cubic polynomial fit
to PSID data:
\begin{equation}
    f(\text{age}) \;=\;
    -2.170 + 0.168\,(\text{age}/10)
    - 0.032\,(\text{age}/10)^2
    + 0.002\,(\text{age}/10)^3.
    \label{eq:income_profile}
\end{equation}
Permanent and transitory income shocks have variances $\sigma_u^2 = 0.0106$
and $\sigma_\varepsilon^2 = 0.0738$ (CGM Table~I).
The retirement replacement rate $\lambda = 0.65$ matches the median
Social Security replacement rate at median lifetime earnings.

Financial asset parameters follow the same source:
the equity premium is $4\%$, stock return volatility $15.7\%$,
and the risk-free rate $2\%$ real.

\subsection{Housing Return Decomposition}
\label{sec:calibration:housing}

Single-location housing return volatility $\sigma_h = 11.5\%$ is from
Cocco (2005).
We decompose the single-location log return into a shared aggregate factor
$\eta_{\text{div}}$ and an idiosyncratic component $\iota_\ell$:
\begin{equation}
    \log R_\ell \;=\; \mu_h + \eta_{\text{div}} + \iota_\ell,
    \qquad
    \eta_{\text{div}} \sim \mathcal{N}(0,\sigma_{\text{div}}^2),
    \quad
    \iota_\ell \sim \mathcal{N}(0,\sigma_\iota^2),
    \label{eq:return_decomp}
\end{equation}
with $\sigma_h^2 = \sigma_{\text{div}}^2 + \sigma_\iota^2$.
We set $\sigma_{\text{div}} = 10\%$ so that the aggregate factor accounts
for approximately $75\%$ of total single-location housing variance
($\sigma_{\text{div}}^2 / \sigma_h^2 \approx 0.756$), consistent with
the common factor dominating MSA-level variance in Case-Shiller data
and with Piazzesi and Schneider~(2016).
The idiosyncratic component $\sigma_\iota = \sqrt{0.115^2 - 0.10^2} \approx 5.7\%$
is calibrated residually.

The idiosyncratic components at the two locations follow a bivariate normal
with correlation $\rho_{AB}$ (Section~\ref{sec:calibration:rhoAB}).
The Cholesky factorisation is:
\begin{align}
    \iota_A &= \sigma_\iota\,\xi_A, \label{eq:chol_A} \\
    \iota_B &= \rho_{AB}\,\iota_A + \sqrt{1 - \rho_{AB}^2}\,\sigma_\iota\,\xi_B,
    \label{eq:chol_B}
\end{align}
with $\xi_A, \xi_B \sim \mathcal{N}(0,1)$ independent.

The rent-to-price ratio $\varrho = 5\%$ follows Yao and Zhang~(2005).
Maintenance costs $m = 1\%$ follow Cocco~(2005).
The rent-saving wedge $\delta_{\text{own}} = \varrho - m = 4\%$ governs
the per-period benefit of occupying a fractional token position
($\kappa$ rule, equation~(\ref{eq:kappa}) in Section~\ref{sec:model:kappa}).

\subsection{Mobility Rate}
\label{sec:calibration:mobility}

We anchor the relocation probability to PSID inter-MSA mobility data.
Total annual relocation rates in the PSID average $7$--$12\%$ for
working-age households but include short-distance, within-MSA moves that
do not change the relevant housing market.
Long-distance, inter-MSA moves account for approximately $40$--$50\%$
of total PSID mobility, yielding an inter-MSA rate of $5$--$7\%$ annually
for prime-age workers.
We set $p_{\text{work}} = 6\%$ (PSID midpoint) for working ages
$25$--$64$ and $p_{\text{ret}} = 2\%$ for ages $65$--$80$.

This choice is consistent with Bagliano, Fugazza, and Nicodano~(2014, RFS),
who calibrate career-driven mobility shocks at $5$--$8\%$ annually, and
broadly in the range of Yao and Zhang~(2005) ($4\%$) and Saks and
Wozniak~(2011, JLE) ($2$--$4\%$ inter-state, implying $4$--$8\%$ inter-MSA).

The sensitivity grid (Table~\ref{tab:sensitivity_grid} below)
spans $p_{\text{work}} \in \{0, 0.02, 0.06, 0.12\}$.
At $p_{\text{work}} = 0$, the cross-location hedge channel must vanish by
construction (no relocation, no cost-saving motive); this is a pre-registered
falsification test (Referee Round~4, test~r).

\subsection{Transaction Costs}
\label{sec:calibration:txcost}

\paragraph{Traditional housing ($E1\_2L$).}
We set the seller commission and closing costs at $\tau_{\text{sell}} = 6\%$,
consistent with the NAR~(2023) aggregate for a typical single-family
home sale ($5$--$6\%$ agent commission plus $0.5$--$1\%$ closing costs).
The buyer cost is $\tau_{\text{buy}} = 2.5\%$,
covering origination fees ($0.75\%$), title insurance and escrow
($0.75\%$), appraisal and inspection ($0.25\%$), and prepaid items
($0.25\%$), consistent with CFPB and ALTA surveys.
The implied round-trip cost for an $E1\_2L$ relocation is
$\tau_{\text{sell}} + \tau_{\text{buy}} = 8.5\%$,
within the $5$--$10\%$ range used in the structural housing literature
(Diaz and Luengo-Prado, 2008; Head, Lloyd-Ellis, and Sun, 2014).

\paragraph{Tokenized housing ($E2\_2L$).}
Platform surveys of RealT, Lofty, and DigiShares~(2024) report trading
fees in the range $0.1$--$0.5\%$ plus blockchain gas and settlement
of $0.05$--$0.3\%$ on Layer-2 networks.
We set $\tau_{\text{token}} = 1\%$ to include a regulatory-compliance
margin of $0.5\%$ (reflecting SEC Regulation~D and AML overhead).
The structural distinction from $E1\_2L$ is that $\tau_{\text{token}} \ll \tau_{\text{sell}} + \tau_{\text{buy}}$:
tokens are portable across relocations, so the $8.5\%$ round-trip is avoided
for every unit of cross-location position retained.

\paragraph{Expected hedge premium.}
The per-period value of pre-positioning one unit of $x_B$ while at
location $A$ is
\begin{equation}
    \text{hedge premium per unit}
    \;=\;
    p_{\text{work}} \times \tau_{\text{buy}}
    \;\approx\;
    0.06 \times 0.025
    \;=\; 0.15\%\ \text{per period}.
    \label{eq:hedge_premium}
\end{equation}
This is the incremental saving on buying costs that accrues when the
household has pre-positioned $x_B > 0$ prior to relocation.
Accumulated over a working-life horizon, equation~(\ref{eq:hedge_premium})
implies a lifetime CEV contribution of $1$--$2\%$ from the hedge
channel alone, with the continuous-$x$ rent-saving channel adding a further
$3$--$4\%$ (see Section~\ref{sec:results:decomp}).

\subsection{Cross-Location House-Price Correlation}
\label{sec:calibration:rhoAB}

The parameter $\rho_{AB}$ governs the idiosyncratic correlation between
location $A$ and location $B$ returns.
Using the decomposition~(\ref{eq:return_decomp}),
the observed return correlation between two MSAs is:
\begin{equation}
    \operatorname{Corr}(R_A, R_B)
    \;\approx\;
    \frac{\sigma_{\text{div}}^2}{\sigma_h^2}
    \;+\;
    \frac{\sigma_\iota^2}{\sigma_h^2}\,\rho_{AB}
    \;\approx\; 0.756 + 0.244\,\rho_{AB}.
    \label{eq:rhoAB_mapping}
\end{equation}
We set $\rho_{AB} = 0.50$, which maps to an observed return correlation of
$\approx 0.88$, consistent with pairs of geographically proximate
US metros in the Case-Shiller 20-city panel (2000--2023).
This is the midpoint of the $0.30$--$0.70$ range that covers inter-regional
moves (different metro areas): at one extreme, Chicago--Phoenix
($\hat{\rho}_{AB} \approx 0.35$) and at the other,
NY--Boston ($\hat{\rho}_{AB} \approx 0.60$).

Equation~(\ref{eq:rhoAB_mapping}) implies that the hedge value of
$x_B$ at location $A$ derives from bearing the idiosyncratic component
$\iota_B$ independently of $\iota_A$.
As $\rho_{AB} \to 1$, locations $A$ and $B$ become nearly identical
markets and $x_B$ is a near-substitute for $x_A$; the hedge channel
collapses.
The Round-4 referee pre-registered test~(m) requires $\operatorname{CEV}(E2\_2L \text{ vs } E1\_2L)$
to decrease significantly at $\rho_{AB} = 0.95$.

\subsection{Sensitivity Grid}
\label{sec:calibration:sensitivity}

Table~\ref{tab:sensitivity_grid} lists the pre-registered sensitivity
sweeps queued for execution after the v4 baseline results are available.
The three P1-priority sweeps cover the three parameters most directly
linked to the welfare decomposition channels.

\begin{table}[t]
\centering
\caption{Pre-Registered Sensitivity Sweeps}
\label{tab:sensitivity_grid}
\small
\begin{tabular}{llll}
\toprule
Parameter & Sweep values & Channel & Script \\
\midrule
$\rho_{AB}$          & $\{0,\, 0.25,\, 0.50,\, 0.75,\, 0.95\}$  & Hedge (maintained) & \texttt{sweep\_rhoAB.sh} \\
$p_{\text{work}}$    & $\{0,\, 0.02,\, 0.06,\, 0.12\}$           & Hedge (frequency)  & \texttt{sweep\_prelocate.sh} \\
$\tau_{\text{buy}}$  & $\{0,\, 0.025,\, 0.04,\, 0.06\}$          & Avoided tx-cost    & \texttt{sweep\_txcost.sh} \\
\midrule
\multicolumn{4}{l}{\textit{Phase 2 sweeps (after H2' approval)}} \\
$\gamma$             & $\{3,\, 5,\, 8\}$                         & Risk channel       & --- \\
$\delta_{\text{own}}$& $\{0.02,\, 0.04,\, 0.06\}$               & Rent-saving        & --- \\
$\sigma_{\text{div}}$& $\{0.08,\, 0.10,\, 0.12\}$               & Idio.\ risk share  & --- \\
\bottomrule
\end{tabular}
\end{table}

At $p_{\text{work}} = 0$ and at $\tau_{\text{buy}} = 0$,
the hedge premium~(\ref{eq:hedge_premium}) is exactly zero; any
remaining welfare gap between $E2\_2L$ and $E1\_2L$ in these cells
reflects the continuous-$x$ rent-saving channel only and serves
as a falsification bound.
At $\rho_{AB} = 0.95$, the maintained-hedge channel must also collapse
because $x_B$ and $x_A$ are near-substitutes.
These three boundary conditions define sharp predictions that distinguish
the two decomposed channels (Section~\ref{sec:results:decomp}).

%              Requires amsmath, booktabs, multirow, siunitx (for tabular).

% ─────────────────────────────────────────────────────────────────────────────
\section{Calibration}
\label{sec:calibration}
% ─────────────────────────────────────────────────────────────────────────────

Table~\ref{tab:parameters} summarises the baseline parameter values.
The calibration follows the household-finance lifecycle tradition
(Cocco, Gomes, and Maenhout, 2005; Yao and Zhang, 2005), extended
with three v4-specific objects: the relocation probability $p_{\text{work}}$,
the round-trip transaction-cost pair $(\tau_{\text{sell}}, \tau_{\text{buy}})$,
and the cross-location idiosyncratic correlation $\rho_{AB}$.
We discuss each group in turn.

% ── Table 1: Parameter summary ───────────────────────────────────────────────
\begin{table}[t]
\centering
\caption{Baseline Parameters}
\label{tab:parameters}
\small
\begin{tabular}{llcl}
\toprule
Parameter & Symbol & Value & Source \\
\midrule
\multicolumn{4}{l}{\textit{Preferences and lifecycle}} \\
Risk aversion         & $\gamma$         & 5.0  & CGM (2005) \\
Time discount         & $\beta$          & 0.96 & CGM (2005) \\
Entry age / retire / terminal & --- & 25 / 65 / 80 & Standard \\
\addlinespace
\multicolumn{4}{l}{\textit{Financial assets}} \\
Risk-free gross return     & $R_f$           & 1.02  & CGM (2005) \\
Equity risk premium        & $\mu_S - r_f$   & 4\%   & CGM (2005) \\
Stock return volatility    & $\sigma_S$      & 15.7\% & CGM (2005) \\
\addlinespace
\multicolumn{4}{l}{\textit{Income process (CGM 2005)}} \\
Permanent shock variance   & $\sigma_u^2$    & 0.0106 & CGM Table~I \\
Transitory shock variance  & $\sigma_\varepsilon^2$ & 0.0738 & CGM Table~I \\
Retirement replacement     & $\lambda$       & 0.65  & CGM (2005) \\
\addlinespace
\multicolumn{4}{l}{\textit{Housing return (Cocco 2005; Yao-Zhang 2005)}} \\
Expected log price growth  & $g_h$           & 1.6\%  & Cocco (2005) \\
Total housing return std   & $\sigma_h$      & 11.5\% & Cocco (2005) \\
Aggregate factor std       & $\sigma_{\text{div}}$ & 10.0\% & Case-Shiller / NAREIT \\
Idiosyncratic std (derived) & $\sigma_\iota$ & 5.7\%  & $\sqrt{\sigma_h^2 - \sigma_{\text{div}}^2}$ \\
Rent-to-price ratio        & $\varrho$       & 5.0\%  & Yao and Zhang (2005) \\
Maintenance-to-price ratio & $m$             & 1.0\%  & Cocco (2005) \\
Rent-saving wedge (derived) & $\delta_{\text{own}} = \varrho - m$ & 4.0\% & --- \\
\addlinespace
\multicolumn{4}{l}{\textit{Mobility (PSID-anchored)}} \\
Relocation prob, working age & $p_{\text{work}}$ & 6.0\% & PSID inter-MSA \\
Relocation prob, retired     & $p_{\text{ret}}$  & 2.0\% & PSID retired \\
\addlinespace
\multicolumn{4}{l}{\textit{Transaction costs}} \\
Traditional sell cost  & $\tau_{\text{sell}}$ & 6.0\% & NAR (2023) \\
Traditional buy cost   & $\tau_{\text{buy}}$  & 2.5\% & CFPB / ALTA \\
Token transfer cost    & $\tau_{\text{token}}$& 1.0\% & RealT / Lofty surveys \\
\addlinespace
\multicolumn{4}{l}{\textit{Cross-location correlation}} \\
Idiosyncratic corr., MSA pair & $\rho_{AB}$ & 0.50 & Case-Shiller (2000--2023) \\
\addlinespace
\multicolumn{4}{l}{\textit{Mortgage (baseline off)}} \\
LTV constraint & $\overline{\text{ltv}}$ & 0.0 & Baseline no-mortgage \\
\bottomrule
\end{tabular}
\end{table}

\subsection{Preferences and Income Process}
\label{sec:calibration:prefs}

Preferences and the income process are taken directly from
Cocco, Gomes, and Maenhout (2005, CGM hereafter).
Risk aversion $\gamma = 5$ and the discount factor $\beta = 0.96$ are
their baseline values.
The labor income age profile follows the CGM cubic polynomial fit
to PSID data:
\begin{equation}
    f(\text{age}) \;=\;
    -2.170 + 0.168\,(\text{age}/10)
    - 0.032\,(\text{age}/10)^2
    + 0.002\,(\text{age}/10)^3.
    \label{eq:income_profile}
\end{equation}
Permanent and transitory income shocks have variances $\sigma_u^2 = 0.0106$
and $\sigma_\varepsilon^2 = 0.0738$ (CGM Table~I).
The retirement replacement rate $\lambda = 0.65$ matches the median
Social Security replacement rate at median lifetime earnings.

Financial asset parameters follow the same source:
the equity premium is $4\%$, stock return volatility $15.7\%$,
and the risk-free rate $2\%$ real.

\subsection{Housing Return Decomposition}
\label{sec:calibration:housing}

Single-location housing return volatility $\sigma_h = 11.5\%$ is from
Cocco (2005).
We decompose the single-location log return into a shared aggregate factor
$\eta_{\text{div}}$ and an idiosyncratic component $\iota_\ell$:
\begin{equation}
    \log R_\ell \;=\; \mu_h + \eta_{\text{div}} + \iota_\ell,
    \qquad
    \eta_{\text{div}} \sim \mathcal{N}(0,\sigma_{\text{div}}^2),
    \quad
    \iota_\ell \sim \mathcal{N}(0,\sigma_\iota^2),
    \label{eq:return_decomp}
\end{equation}
with $\sigma_h^2 = \sigma_{\text{div}}^2 + \sigma_\iota^2$.
We set $\sigma_{\text{div}} = 10\%$ so that the aggregate factor accounts
for approximately $75\%$ of total single-location housing variance
($\sigma_{\text{div}}^2 / \sigma_h^2 \approx 0.756$), consistent with
the common factor dominating MSA-level variance in Case-Shiller data
and with Piazzesi and Schneider~(2016).
The idiosyncratic component $\sigma_\iota = \sqrt{0.115^2 - 0.10^2} \approx 5.7\%$
is calibrated residually.

The idiosyncratic components at the two locations follow a bivariate normal
with correlation $\rho_{AB}$ (Section~\ref{sec:calibration:rhoAB}).
The Cholesky factorisation is:
\begin{align}
    \iota_A &= \sigma_\iota\,\xi_A, \label{eq:chol_A} \\
    \iota_B &= \rho_{AB}\,\iota_A + \sqrt{1 - \rho_{AB}^2}\,\sigma_\iota\,\xi_B,
    \label{eq:chol_B}
\end{align}
with $\xi_A, \xi_B \sim \mathcal{N}(0,1)$ independent.

The rent-to-price ratio $\varrho = 5\%$ follows Yao and Zhang~(2005).
Maintenance costs $m = 1\%$ follow Cocco~(2005).
The rent-saving wedge $\delta_{\text{own}} = \varrho - m = 4\%$ governs
the per-period benefit of occupying a fractional token position
($\kappa$ rule, equation~(\ref{eq:kappa}) in Section~\ref{sec:model:kappa}).

\subsection{Mobility Rate}
\label{sec:calibration:mobility}

We anchor the relocation probability to PSID inter-MSA mobility data.
Total annual relocation rates in the PSID average $7$--$12\%$ for
working-age households but include short-distance, within-MSA moves that
do not change the relevant housing market.
Long-distance, inter-MSA moves account for approximately $40$--$50\%$
of total PSID mobility, yielding an inter-MSA rate of $5$--$7\%$ annually
for prime-age workers.
We set $p_{\text{work}} = 6\%$ (PSID midpoint) for working ages
$25$--$64$ and $p_{\text{ret}} = 2\%$ for ages $65$--$80$.

This choice is consistent with Bagliano, Fugazza, and Nicodano~(2014, RFS),
who calibrate career-driven mobility shocks at $5$--$8\%$ annually, and
broadly in the range of Yao and Zhang~(2005) ($4\%$) and Saks and
Wozniak~(2011, JLE) ($2$--$4\%$ inter-state, implying $4$--$8\%$ inter-MSA).

The sensitivity grid (Table~\ref{tab:sensitivity_grid} below)
spans $p_{\text{work}} \in \{0, 0.02, 0.06, 0.12\}$.
At $p_{\text{work}} = 0$, the cross-location hedge channel must vanish by
construction (no relocation, no cost-saving motive); this is a pre-registered
falsification test (Referee Round~4, test~r).

\subsection{Transaction Costs}
\label{sec:calibration:txcost}

\paragraph{Traditional housing ($E1\_2L$).}
We set the seller commission and closing costs at $\tau_{\text{sell}} = 6\%$,
consistent with the NAR~(2023) aggregate for a typical single-family
home sale ($5$--$6\%$ agent commission plus $0.5$--$1\%$ closing costs).
The buyer cost is $\tau_{\text{buy}} = 2.5\%$,
covering origination fees ($0.75\%$), title insurance and escrow
($0.75\%$), appraisal and inspection ($0.25\%$), and prepaid items
($0.25\%$), consistent with CFPB and ALTA surveys.
The implied round-trip cost for an $E1\_2L$ relocation is
$\tau_{\text{sell}} + \tau_{\text{buy}} = 8.5\%$,
within the $5$--$10\%$ range used in the structural housing literature
(Diaz and Luengo-Prado, 2008; Head, Lloyd-Ellis, and Sun, 2014).

\paragraph{Tokenized housing ($E2\_2L$).}
Platform surveys of RealT, Lofty, and DigiShares~(2024) report trading
fees in the range $0.1$--$0.5\%$ plus blockchain gas and settlement
of $0.05$--$0.3\%$ on Layer-2 networks.
We set $\tau_{\text{token}} = 1\%$ to include a regulatory-compliance
margin of $0.5\%$ (reflecting SEC Regulation~D and AML overhead).
The structural distinction from $E1\_2L$ is that $\tau_{\text{token}} \ll \tau_{\text{sell}} + \tau_{\text{buy}}$:
tokens are portable across relocations, so the $8.5\%$ round-trip is avoided
for every unit of cross-location position retained.

\paragraph{Expected hedge premium.}
The per-period value of pre-positioning one unit of $x_B$ while at
location $A$ is
\begin{equation}
    \text{hedge premium per unit}
    \;=\;
    p_{\text{work}} \times \tau_{\text{buy}}
    \;\approx\;
    0.06 \times 0.025
    \;=\; 0.15\%\ \text{per period}.
    \label{eq:hedge_premium}
\end{equation}
This is the incremental saving on buying costs that accrues when the
household has pre-positioned $x_B > 0$ prior to relocation.
Accumulated over a working-life horizon, equation~(\ref{eq:hedge_premium})
implies a lifetime CEV contribution of $1$--$2\%$ from the hedge
channel alone, with the continuous-$x$ rent-saving channel adding a further
$3$--$4\%$ (see Section~\ref{sec:results:decomp}).

\subsection{Cross-Location House-Price Correlation}
\label{sec:calibration:rhoAB}

The parameter $\rho_{AB}$ governs the idiosyncratic correlation between
location $A$ and location $B$ returns.
Using the decomposition~(\ref{eq:return_decomp}),
the observed return correlation between two MSAs is:
\begin{equation}
    \operatorname{Corr}(R_A, R_B)
    \;\approx\;
    \frac{\sigma_{\text{div}}^2}{\sigma_h^2}
    \;+\;
    \frac{\sigma_\iota^2}{\sigma_h^2}\,\rho_{AB}
    \;\approx\; 0.756 + 0.244\,\rho_{AB}.
    \label{eq:rhoAB_mapping}
\end{equation}
We set $\rho_{AB} = 0.50$, which maps to an observed return correlation of
$\approx 0.88$, consistent with pairs of geographically proximate
US metros in the Case-Shiller 20-city panel (2000--2023).
This is the midpoint of the $0.30$--$0.70$ range that covers inter-regional
moves (different metro areas): at one extreme, Chicago--Phoenix
($\hat{\rho}_{AB} \approx 0.35$) and at the other,
NY--Boston ($\hat{\rho}_{AB} \approx 0.60$).

Equation~(\ref{eq:rhoAB_mapping}) implies that the hedge value of
$x_B$ at location $A$ derives from bearing the idiosyncratic component
$\iota_B$ independently of $\iota_A$.
As $\rho_{AB} \to 1$, locations $A$ and $B$ become nearly identical
markets and $x_B$ is a near-substitute for $x_A$; the hedge channel
collapses.
The Round-4 referee pre-registered test~(m) requires $\operatorname{CEV}(E2\_2L \text{ vs } E1\_2L)$
to decrease significantly at $\rho_{AB} = 0.95$.

\subsection{Sensitivity Grid}
\label{sec:calibration:sensitivity}

Table~\ref{tab:sensitivity_grid} lists the pre-registered sensitivity
sweeps queued for execution after the v4 baseline results are available.
The three P1-priority sweeps cover the three parameters most directly
linked to the welfare decomposition channels.

\begin{table}[t]
\centering
\caption{Pre-Registered Sensitivity Sweeps}
\label{tab:sensitivity_grid}
\small
\begin{tabular}{llll}
\toprule
Parameter & Sweep values & Channel & Script \\
\midrule
$\rho_{AB}$          & $\{0,\, 0.25,\, 0.50,\, 0.75,\, 0.95\}$  & Hedge (maintained) & \texttt{sweep\_rhoAB.sh} \\
$p_{\text{work}}$    & $\{0,\, 0.02,\, 0.06,\, 0.12\}$           & Hedge (frequency)  & \texttt{sweep\_prelocate.sh} \\
$\tau_{\text{buy}}$  & $\{0,\, 0.025,\, 0.04,\, 0.06\}$          & Avoided tx-cost    & \texttt{sweep\_txcost.sh} \\
\midrule
\multicolumn{4}{l}{\textit{Phase 2 sweeps (after H2' approval)}} \\
$\gamma$             & $\{3,\, 5,\, 8\}$                         & Risk channel       & --- \\
$\delta_{\text{own}}$& $\{0.02,\, 0.04,\, 0.06\}$               & Rent-saving        & --- \\
$\sigma_{\text{div}}$& $\{0.08,\, 0.10,\, 0.12\}$               & Idio.\ risk share  & --- \\
\bottomrule
\end{tabular}
\end{table}

At $p_{\text{work}} = 0$ and at $\tau_{\text{buy}} = 0$,
the hedge premium~(\ref{eq:hedge_premium}) is exactly zero; any
remaining welfare gap between $E2\_2L$ and $E1\_2L$ in these cells
reflects the continuous-$x$ rent-saving channel only and serves
as a falsification bound.
At $\rho_{AB} = 0.95$, the maintained-hedge channel must also collapse
because $x_B$ and $x_A$ are near-substitutes.
These three boundary conditions define sharp predictions that distinguish
the two decomposed channels (Section~\ref{sec:results:decomp}).
